\chapter{Introducción}

Este capítulo introduce los antecedentes, objetivos y plan de
trabajo de la memoria, tal y como exige la normativa del trabajo
de fin de estudios de la Facultad de Informática
(\href{https://informatica.ucm.es/normativa-tfm-2425}{informatica.ucm.es/normativa-tfm-2425}).

\section{Normativa}

\begin{itemize}\setlength{\itemsep}{0pt}
  \item Extensión mínima: 25 páginas para TFG de un solo autor (más
        5 por cada coautor adicional); 50 páginas para TFM.
  \item Máximo 10 palabras clave en cada idioma.
  \item El resumen en inglés del TFM debe ocupar
        aproximadamente media página.
  \item Todo material no original (texto, figuras, código) debe
        citarse y respetar licencias.
  \item \textbf{Sólo TFM}: el estudiante debe firmar y adjuntar la
        \emph{declaración responsable sobre autoría y uso ético de
        herramientas de inteligencia artificial}.
\end{itemize}

\section{Antecedentes}

La Facultad de Informática de la UCM ha utilizado durante más de
una década el sistema TeXiS~\cite{texis} para sus memorias de
tesis y trabajos de fin de carrera. Esta cáscara ayuda a preparar
documentos complejos con buenas prácticas, pero al ser una
cáscara y no una plantilla, es más complicada de integrar con
otros sistemas, además de no incentivar el aprendizaje de LaTeX
por los usuarios.

\section{Objetivos}

Los objetivos de este trabajo son:
\begin{itemize}
  \item Disponer de una plantilla moderna basada en LuaLaTeX.
  \item Separar los aspectos normativos de los estéticos.
  \item Ofrecer una base mínima que no interfiera con los
        paquetes que cada estudiante decida usar.
\end{itemize}
